\slajdid{02-s032b}
\begin{frame}[label=02-s032b]{Minimizacija logičkih funkcija — primer}
\setlength{\abovedisplayskip}{4pt}\setlength{\abovedisplayshortskip}{4pt}
\setlength{\belowdisplayskip}{4pt}\setlength{\belowdisplayshortskip}{4pt}
% element: 02-s032b:e01
Indeksi dva proizvoda su 4 i 6, odnosno \(100_2\) i \(110_2\). Hemingovo rastojanje, broj različitih bitova, jeste jedan: razlikuju se samo u drugom bitu, odnosno u literalima \(\bar B\) i \(B\). Zato je moguće:
\[F=\cdots+C\bar B\bar A+CB\bar A+\cdots
 =\cdots+C\bar A(\bar B+B)+\cdots
 =\cdots+C\bar A+\cdots.\]
Umesto dva proizvoda, za koja su bila potrebna dva troulazna I kola, dovoljan je jedan proizvod i jedno dvoulazno I kolo. Na završno ILI kolo dovodi se izlaz jednog I kola umesto dva.
\par\vspace{3mm}
% element: 02-s032b:e02
\Rezultat{Broj I kola smanjujemo sa dva na jedno, a \alert{ukupan broj ulaza sa \(2\cdot3+2=8\) na \(1\cdot2+1=3\)}.}
\par\vspace{3mm}
% element: 02-s032b:e03
U CMOS logici time \alert{značajno smanjujemo broj potrebnih tranzistora, odnosno površinu} koju zauzima taj deo kombinacione mreže.
\end{frame}
